
\documentclass[a4paper,11pt]{article}
\usepackage[a4paper, margin=8em]{geometry}

% usa i pacchetti per la scrittura in italiano
\usepackage[french,italian]{babel}
\usepackage[T1]{fontenc}
\usepackage[utf8]{inputenc}
\frenchspacing 

% usa i pacchetti per la formattazione matematica
\usepackage{amsmath, amssymb, amsthm, amsfonts}

% usa altri pacchetti
\usepackage{gensymb}
\usepackage{hyperref}
\usepackage{standalone}

% imposta il titolo
\title{Appunti Comunicazioni Numeriche}
\author{Luca Seggiani}
\date{2026}

% imposta lo stile
% usa helvetica
\usepackage[scaled]{helvet}
% usa palatino
\usepackage{palatino}
% usa un font monospazio guardabile
\usepackage{lmodern}

\renewcommand{\rmdefault}{ppl}
\renewcommand{\sfdefault}{phv}
\renewcommand{\ttdefault}{lmtt}

% disponi il titolo
\makeatletter
\renewcommand{\maketitle} {
	\begin{center} 
		\begin{minipage}[t]{.8\textwidth}
			\textsf{\huge\bfseries \@title} 
		\end{minipage}%
		\begin{minipage}[t]{.2\textwidth}
			\raggedleft \vspace{-1.65em}
			\textsf{\small \@author} \vfill
			\textsf{\small \@date}
		\end{minipage}
		\par
	\end{center}

	\thispagestyle{empty}
	\pagestyle{fancy}
}
\makeatother

% disponi teoremi
\usepackage{tcolorbox}
\newtcolorbox[auto counter, number within=section]{theorem}[2][]{%
	colback=blue!10, 
	colframe=blue!40!black, 
	sharp corners=northwest,
	fonttitle=\sffamily\bfseries, 
	title=~\thetcbcounter: #2, 
	#1
}

% disponi definizioni
\newtcolorbox[auto counter, number within=section]{definition}[2][]{%
	colback=red!10,
	colframe=red!40!black,
	sharp corners=northwest,
	fonttitle=\sffamily\bfseries,
	title=~\thetcbcounter: #2,
	#1
}

% U.D.M
\newcommand{\amp}{\ensuremath{\mathrm{A}}}
\newcommand{\volt}{\ensuremath{\mathrm{V}}}
\newcommand{\meter}{\ensuremath{\mathrm{m}}}
\newcommand{\second}{\ensuremath{\mathrm{s}}}
\newcommand{\farad}{\ensuremath{\mathrm{F}}}
\newcommand{\henry}{\ensuremath{\mathrm{H}}}
\newcommand{\siemens}{\ensuremath{\mathrm{S}}}

% circuiti
\usepackage{circuitikz}
\usetikzlibrary{babel}

% disegni
\usepackage{pgfplots}
\pgfplotsset{width=10cm,compat=1.9}

% disponi codice
\usepackage{listings}
\usepackage[table]{xcolor}

\lstdefinestyle{codestyle}{
		backgroundcolor=\color{black!5}, 
		commentstyle=\color{codegreen},
		keywordstyle=\bfseries\color{magenta},
		numberstyle=\sffamily\tiny\color{black!60},
		stringstyle=\color{green!50!black},
		basicstyle=\ttfamily\footnotesize,
		breakatwhitespace=false,         
		breaklines=true,                 
		captionpos=b,                    
		keepspaces=true,                 
		numbers=left,                    
		numbersep=5pt,                  
		showspaces=false,                
		showstringspaces=false,
		showtabs=false,                  
		tabsize=2
}

\lstdefinestyle{shellstyle}{
		backgroundcolor=\color{black!5}, 
		basicstyle=\ttfamily\footnotesize\color{black}, 
		commentstyle=\color{black}, 
		keywordstyle=\color{black},
		numberstyle=\color{black!5},
		stringstyle=\color{black}, 
		showspaces=false,
		showstringspaces=false, 
		showtabs=false, 
		tabsize=2, 
		numbers=none, 
		breaklines=true
}

\lstdefinelanguage{javascript}{
	keywords={typeof, new, true, false, catch, function, return, null, catch, switch, var, if, in, while, do, else, case, break},
	keywordstyle=\color{blue}\bfseries,
	ndkeywords={class, export, boolean, throw, implements, import, this},
	ndkeywordstyle=\color{darkgray}\bfseries,
	identifierstyle=\color{black},
	sensitive=false,
	comment=[l]{//},
	morecomment=[s]{/*}{*/},
	commentstyle=\color{purple}\ttfamily,
	stringstyle=\color{red}\ttfamily,
	morestring=[b]',
	morestring=[b]"
}

% disponi sezioni
\usepackage{titlesec}

\titleformat{\section}
	{\sffamily\Large\bfseries} 
	{\thesection}{1em}{} 
\titleformat{\subsection}
	{\sffamily\large\bfseries}   
	{\thesubsection}{1em}{} 
\titleformat{\subsubsection}
	{\sffamily\normalsize\bfseries} 
	{\thesubsubsection}{1em}{}

% disponi alberi
\usepackage{forest}

\forestset{
	rectstyle/.style={
		for tree={rectangle,draw,font=\large\sffamily}
	},
	roundstyle/.style={
		for tree={circle,draw,font=\large}
	}
}

% disponi algoritmi
\usepackage{algorithm}
\usepackage{algorithmic}
\makeatletter
\renewcommand{\ALG@name}{Algoritmo}
\makeatother

% disponi numeri di pagina
\usepackage{fancyhdr}
\fancyhf{} 
\fancyfoot[L]{\sffamily{\thepage}}

\makeatletter
\fancyhead[L]{\raisebox{1ex}[0pt][0pt]{\sffamily{\@title \ \@date}}} 
\fancyhead[R]{\raisebox{1ex}[0pt][0pt]{\sffamily{\@author}}}
\makeatother

\begin{document}
% sezione (data)
\section{Lezione del 22-05-26}

% stili pagina
\thispagestyle{empty}
\pagestyle{fancy}

% testo
\subsubsection{Dimensionamento dei filtri sull'SNR}
Nella scorsa lezione abbiamo visto il dimensionamento dei filtri di ricezione e di trasmissione, e avevamo visto che per rimuovere l'ISI, dato:
$$
g(t) = g_R(t) \circledast g_T(t)
$$
erano condizioni per la banda $B_G$:
$$
2B_G \leq \frac{1}{T_S} \ (\text{Nyquist})
, \quad
B_G \geq \frac{1}{2 T_S} \ (\text{ISI})
$$
di cui una condizione di Nyquist, e l'altra sostanzialmente equivalente alla condizione di Nyquist, ma all'inverso (vogliamo esplicitamente l'overlap degli spettri immagine per l'annullamento dell'ISI).

Ci chiediamo quindi come dimensionare i due filtri separatamente. Occorrerà introdurre un'ulteriore condizione, che troveremo nella massimizzazione dell'\textit{SNR}. Preso il simbolo al ricevitore:
$$
x(k) = \sqrt{\beta} a_k g(0) + n(k) \implies y(k) = \frac{x(k)}{\sqrt{\beta} g(0)} = a_l + \frac{n(k)}{\sqrt{\beta} g(0)}
$$

Vorremo calcolare il rapporto \textit{segnale-rumore} (\textbf{SNR}, \textit{Signal to Noise Ratio}) \textit{medio} fra il simbolo $a_k$ e il rumore in $n(k)$:
$$
\mathrm{SNR} = \frac{E \{ |a_k|^2 \}}{ E \{ |n(k)|^2 \} } \sqrt{\beta} g(0) = \underbrace{\frac{M^2 - 1}{3} \frac{2\beta}{N_0}}_\text{costante $C$} \frac{g^2(0)}{E_{g_R}}
$$
dato il fatto che la costellazione ha media nulla, il rumore ha media nulla, e abbiamo già calcolato le varianze di entrambi ($\frac{M^2 - 1}{3}$ per il segnale e la potenza $P_n$ per il rumore). Isoliamo inoltre la costante data dai parametri del sistema (e non dai filtri) a sinistra.
Esplicitando i termini che ci servono abbiamo:
$$
\mathrm{SNR} = C \frac{ \left( \int_{-\infty}^{\infty} g_T(\alpha) g_R(t - \alpha) \, d\alpha \big|_{t = 0} \right)^2 }{ \int_{-\infty}^{\infty} g_R^2 (\alpha) \, d\alpha }
$$

Sfruttiamo la diseguaglianza di \textit{Cauchy-Schwarz} per dire:
$$
\left( \int_{-\infty}^{\infty} g_T(\alpha) g_R(t - \alpha) \, d\alpha \Big|_{t = 0} \right)^2 \leq \left( \int_{-\infty}^{\infty} g_T(\alpha) \, d\alpha \right)^2 \left( \int_{-\infty}^{\infty} g_R(t - \alpha) \, d\alpha \Big|_{t = 0} \right)^2
$$
Notiamo che la diseguglianza diventa un uguaglianza quando i due filtri sono linearmente dipendenti, cioè esiste una costante di proporzionalità $\gamma$ per cui:
$$
g_T(\alpha) = \gamma g_R(t - \alpha) \Big|_{t = 0}
$$

Abbiamo allora ottenuto:
$$
\max \left( \mathrm{SNR} \right) = C \frac{ \left( \int_{-\infty}^{\infty} \gamma g_R(-t) \, dt \right)^2 \left( \int_{-\infty}^{\infty} g_R(-\alpha) \, d\alpha \right)^2 }{ \int_{-\infty}^{\infty} g_R^2 (\alpha) \, d\alpha } = C \int_{-\infty}^{\infty} \gamma g_R(-t) \, dt
$$

Il filtro $g_T(t)$ così formulato:
$$
g_T(t) = \gamma g_R(-t)
$$
viene detto \textit{filtro adattato} al ricevitore, sul trasmettitore.

Abbiamo quindi le 2 condizioni:
$$
\begin{cases}
g_R(t) = g_T(-t) \\
\begin{cases}
g(m T_s) = 0, \quad \forall m \neq 0 \\
g(m T_s) = g(0), \quad m = 0
\end{cases}
\end{cases}
$$
In frequenza la prima condizione viene posta come:
$$
G_R(f) = G_T^*(f), \quad G(f) = G_T(f) G_R(f)
$$
per cui:
$$
G(f) = G_T(f) G_T^*(f) = |G(f)|^2 \implies
|G_R(f)| = |G_T(f)| = \sqrt{G(f)}
$$

\subsubsection{Funzioni di filtraggio tipiche}

Facciamo quindi un dimensionamento vero e proprio. Il filtro più semplice da prendere sarà la $\mathrm{rect}$ in tempo: 
$$
g_T(t) = \mathrm{rect} \left( \frac{t}{T_S} \right)
$$
per cui la $g_R(t)$ in tempo sarà uguale a:
$$
g_R(t) = g_T(-t)
$$
La convoluzione fra i due filtri in dominio tempo sarà data da:
$$
g_T(t) \circledast g_T(-t) = \mathrm{tri} \left( \frac{t}{2 T_S} \right)
$$
che è un impulso di Nyquist con valore nullo agli estremi $m T_S$ già da $m = \pm 1$ (annulla l'ISI), e che massimizza l'SNR.

Il problema di usare una $\mathrm{rect}$ (sostanzialmente l'interpolatore a mantenimento in uscita dal trasmettitore in banda base) è che si ha una grande quantità di emissioni fuori banda, date dalla densità spettrale di potenza di $\mathrm{sinc}^2$. Cerchiamo allora di usare l'interpolatore cardinale: il problema sarà dato dal fatto che non possiamo implementarlo nella realtà senza introdurre ritardi.

La soluzione è data dall'interpolatore \textbf{coseno rialzato}:
$$
H(f)=
\begin{cases}
T, & |f|\le \frac{1-\beta}{2T} \\[6pt]
\frac{T}{2}\left[1+\cos\left(\frac{\pi T}{\beta}\left(|f|-\frac{1-\beta}{2T}\right)\right)\right], 
& \frac{1-\beta}{2T}<|f|\le \frac{1+\beta}{2T} \\[10pt]
0, & |f|>\frac{1+\beta}{2T}
\end{cases}
$$
dove $\beta \in [0, 1]$ viene detto \textit{fattore di Rolloff}. Notiamo che il coseno rialzato inizia a decadere a $\frac{1 - \alpha}{2T}$, e finisce a $\frac{1 + \alpha}{2 T}$, con punto medio del \textit{rolloff} (o decadimento) in $\frac{1}{2T}$. La banda del filtro sarà $\frac{1 + \alpha}{2T}$, con valori di $\alpha$ intorno a $\sim 0.3$.

\end{document}
